\documentclass[11pt]{article}
\usepackage[T1]{fontenc}
\usepackage[utf8]{inputenc}
\usepackage[spanish]{babel}
\usepackage{amsmath,amssymb}
\usepackage[a4paper,margin=2.3cm]{geometry}
\usepackage{graphicx}
\usepackage{tikz}
\usetikzlibrary{arrows.meta,babel}
\spanishdecimal{.}
\setlength{\parindent}{0pt}
\setlength{\parskip}{0.35em}
\setlength{\abovedisplayskip}{4pt}
\setlength{\belowdisplayskip}{4pt}
\setlength{\abovedisplayshortskip}{4pt}
\setlength{\belowdisplayshortskip}{4pt}

\newcommand{\encabezado}[3]{%
  {\Large\bfseries #1\par}
  \vspace{2pt}
  {\normalsize #2\hfill #3\par}
  \vspace{6pt}\hrule\vspace{8pt}}

\newcommand{\problema}[2]{%
  \vspace{0.55em}
  \noindent\textbf{#1.}\enspace{\small #2}\par
  \vspace{0.3em}}


% Figuras en línea, junto al paso del desarrollo que ilustran.
\newsavebox{\EDdibujo}
\makeatletter
\newenvironment{EDfigura}{%
  \par\medskip\noindent
  \begin{minipage}{\linewidth}
  \def\@captype{figure}%
  \centering
  \begin{lrbox}{\EDdibujo}%
}{%
  \end{minipage}\par\medskip
}
\newcommand{\EDpie}[2]{%
  \end{lrbox}%
  \ifdim\wd\EDdibujo>\linewidth
    \resizebox{\linewidth}{!}{\usebox{\EDdibujo}}%
  \else
    \usebox{\EDdibujo}%
  \fi
  \caption{#2}\label{#1}%
}
\makeatother

\begin{document}



\encabezado{Inducción electromagnética y ecuaciones de Maxwell}
{Electrodinámica I}{P.S., J.A., J.R.}

\section*{Enunciados de los Problemas}



\subsection*{Problema 1: Fuerza Electromotriz de Movimiento vs. Ley del Flujo}
Un alambre rectilíneo muy largo se ubica sobre el eje $z$ y transporta una corriente estacionaria $I$ en la dirección $+\hat{\mathbf{z}}$. En el semiplano $\phi = 0$ (coordenadas cilíndricas $(s, \phi, z)$), una espira conductora rectangular de ancho radial $a$, altura $b$ (paralela al eje $z$) y resistencia eléctrica total $R$ se aleja perpendicularmente del alambre con velocidad constante $\mathbf{v} = v_0 \hat{\mathbf{s}}$. En un instante arbitrario $t$, el borde vertical más cercano de la espira se encuentra a una distancia $s(t)$ del alambre.
\begin{enumerate}
    \item[a)] Calcule el flujo magnético $\Phi_B(t)$ que atraviesa la espira y utilice la regla del flujo de Faraday ($\mathcal{E} = -d\Phi_B/dt$) para determinar la magnitud y el sentido de circulación de la corriente inducida $I_{\text{ind}}$.
    \item[b)] Vuelva a calcular la fuerza electromotriz $\mathcal{E}$ integrando explícitamente la fuerza magnética por unidad de carga a lo largo del contorno de la espira: $\mathcal{E} = \oint (\mathbf{v} \times \mathbf{B}) \cdot d\mathbf{l}$. Identifique qué segmentos del rectángulo contribuyen a la fem y compruebe que el resultado coincide con el inciso anterior.
    \item[c)] Determine la fuerza magnética neta $\mathbf{F}_{\text{mag}}$ que ejerce el campo del alambre sobre la espira con corriente inducida y verifique que se opone al movimiento (Ley de Lenz). Calcule la potencia mecánica externa $P_{\text{mec}} = \mathbf{F}_{\text{ext}} \cdot \mathbf{v}$ requerida para mantener la velocidad constante y demuestre que es idéntica a la potencia disipada por efecto Joule ($I_{\text{ind}}^2 R$).
\end{enumerate}

\subsection*{Problema 2: Campo Eléctrico Inducido No Conservativo}
Un solenoide ideal muy largo de radio $R$ y $n$ espiras por unidad de longitud tiene su eje coincidente con el eje $z$ y transporta una corriente que crece linealmente con el tiempo según $I(t) = I_0 + \alpha t$, donde $I_0$ y $\alpha$ son constantes positivas.
\begin{enumerate}
    \item[a)] Determine el campo magnético $\mathbf{B}(t)$ en el interior ($s < R$) y en el exterior ($s > R$) del solenoide bajo la aproximación cuasiestática.
    \item[b)] Utilizando la Ley de Faraday en su forma integral y argumentos de simetría cilíndrica, encuentre el vector campo eléctrico inducido $\mathbf{E}(s)$ tanto en el interior ($s < R$) como en el exterior ($s > R$) del solenoide. Compruebe explícitamente que $\nabla \times \mathbf{E} = -\partial \mathbf{B}/\partial t$ en ambas regiones.
    \item[c)] Considere los puntos $P_1 = (2R, 0, 0)$ y $P_2 = (-2R, 0, 0)$ situados fuera del solenoide. Calcule la integral de línea $-\int_{P_1}^{P_2} \mathbf{E} \cdot d\mathbf{l}$ a lo largo de dos caminos semicirculares distintos de radio $2R$ contenidos en el plano $xy$: el Camino 1 por el semiplano superior ($y > 0$) y el Camino 2 por el semiplano inferior ($y < 0$). Explique físicamente por qué ambos resultados difieren a pesar de evaluarse en una región donde $\mathbf{B} = \mathbf{0}$, y justifique por qué no es posible definir un potencial electrostático escalar unívoco $V$ tal que $\mathbf{E} = -\nabla V$.
\end{enumerate}
\newpage
\subsection*{Problema 3: Inductancia Mutua y Reciprocidad de Neumann}
Considere un alambre conductor infinito doblado en ángulo recto sobre los semiejes positivos $x$ e $y$: la corriente ingresa desde $x = +\infty$ hacia el origen a lo largo del semieje $+x$, y sale desde el origen hacia $y = +\infty$ a lo largo del semieje $+y$. En el primer cuadrante del plano $xy$ se coloca una pequeña espira circular plana de radio $a$, cuyo centro se ubica en el punto $(d, d, 0)$, a gran distancia de ambos ejes ($d \gg a$).
\begin{enumerate}
    \item[a)] Argumente brevemente por qué calcular el flujo magnético generado por la pequeña espira circular a través del circuito infinito en ángulo recto es analíticamente complejo, mientras que el cálculo recíproco resulta elemental.
    \item[b)] Asumiendo que por el alambre en ángulo recto circula una corriente estacionaria de prueba $I$, utilice la Ley de Biot-Savart para calcular el vector campo magnético $\mathbf{B}$ exactamente en el centro de la espira circular $(d, d, 0)$.
    \item[c)] Aproximando el campo magnético como uniforme sobre la superficie de la pequeña espira ($a \ll d$), determine el coeficiente de inductancia mutua $M$ entre ambos circuitos.
    \item[d)] Si en la espira circular pequeña se inyecta una corriente alterna $I_c(t) = I_0 \cos(\omega t)$, determine la fuerza electromotriz $\mathcal{E}(t)$ inducida sobre el alambre infinito en ángulo recto.
\end{enumerate}

\subsection*{Problema 4: Autoinductancia Interna y Externa de un Cable Coaxial}
Un cable coaxial rectilíneo muy largo está formado por un conductor cilíndrico interno \textbf{macizo} de radio $a$ y una cáscara conductora cilíndrica externa muy delgada de radio $b$ ($b > a$). Una corriente estacionaria $I$ fluye en dirección $+\hat{\mathbf{z}}$ distribuida uniformemente en toda la sección transversal del núcleo interno, y retorna en dirección $-\hat{\mathbf{z}}$ por la cáscara externa. El medio en todo el sistema tiene permeabilidad $\mu_0$.
\begin{enumerate}
    \item[a)] Determine el vector campo magnético $\mathbf{B}(s)$ en las tres regiones del espacio: $s < a$, $a < s < b$, y $s > b$.
    \item[b)] Calcule la energía magnética almacenada por unidad de longitud $U_B/\ell$ integrando la densidad de energía magnética $u_B = \frac{1}{2\mu_0} B^2$ en todo el espacio transversal.
    \item[c)] A partir de la definición energética $U_B = \frac{1}{2} L I^2$, determine la autoinductancia total por unidad de longitud $\mathcal{L} = L/\ell$, identificando por separado la contribución de la inductancia interna ($s < a$) y de la inductancia externa ($a < s < b$).
    \item[d)] Explique conceptual y matemáticamente por qué si se calcula ingenuamente el flujo magnético $\Phi_B = \int_0^b B(s) \ell \, ds$ y se divide por $I$, se obtiene un factor erróneo para la autoinductancia interna del conductor macizo. Muestre cómo ponderar el enlace de flujo parcial para recuperar el resultado exacto.
\end{enumerate}
\newpage
\subsection*{Problema 5: Ecuación de Ampère-Maxwell en un Alambre Semi-Infinito}
Considere un alambre rectilíneo, muy delgado y semi-infinito que se extiende a lo largo del semieje negativo $z$ (desde $z = -\infty$ hasta terminar en el origen $z = 0$). Por el alambre fluye una corriente constante $I$ en la dirección $+\hat{\mathbf{z}}$. Dado que el conductor termina en el origen, se acumula carga eléctrica puntualmente en $(0,0,0)$ a razón constante $dQ/dt = I$, asumiendo $Q(0) = 0$.
\begin{enumerate}
    \item[a)] Bajo la aproximación cuasiestática, escriba el campo eléctrico $\mathbf{E}(\mathbf{r}, t)$ producido en todo el espacio por la carga puntual $Q(t)$ en el origen (asuma que el alambre permanece eléctricamente neutro). Calcule la densidad de corriente de desplazamiento $\mathbf{J}_d(\mathbf{r})$ para $r > 0$ y compruebe que el flujo neto de la corriente generalizada $\mathbf{J} + \mathbf{J}_d$ a través de una esfera cerrada centrada en el origen es idénticamente cero.
    \item[b)] Sea un punto de observación en coordenadas esféricas $(r, \theta, \phi)$. Defina un lazo circular amperiano $C$ centrado en el eje $z$, de radio cilíndrico $s = r\sin\theta$ y ubicado en el plano horizontal $z = r\cos\theta$. Elija como superficie abierta $S_1$ bordeada por $C$ el casquete esférico de radio constante $r$ que cubre el eje $+z$ ($0 \le \theta' \le \theta$, el cual no es perforado por el alambre). Calcule la corriente de desplazamiento total $I_{d,\text{enc}}$ que atraviesa $S_1$.
    \item[c)] Argumentando por simetría que $\mathbf{B} = B_\phi(r, \theta) \hat{\boldsymbol{\phi}}$, aplique la Ley de Ampère-Maxwell sobre el casquete $S_1$ para hallar $\mathbf{B}(r, \theta)$ en todo el espacio. Verifique que al utilizar el casquete esférico complementario $S_2$ ($\theta \le \theta' \le \pi$, el cual sí es atravesado por el alambre en $\theta' = \pi$), la suma $I_{\text{enc}} + I_{d,\text{enc}}$ arroja exactamente el mismo campo magnético.
    \item[d)] Calcule directamente el campo magnético que generaría solo la corriente de conducción $I$ del alambre semi-infinito mediante la Ley de Biot-Savart. Compare este resultado con el campo total del inciso c), deduzca cuánto vale el campo magnético generado por la distribución radial de corrientes de desplazamiento $\mathbf{J}_d$ y explique físicamente este fenómeno.
\end{enumerate}

\newpage

\section*{Soluciones}

\subsection*{Solución Problema 1}

\paragraph{a) Flujo magnético y corriente inducida:}
Por simetría, el campo del alambre es azimutal. Aplicamos
la ley de Ampère en una circunferencia de radio $s'$ centrada en
el alambre. En esa curva, $B_\phi(s')$ es constante y tangente:
\[
    \oint_{\mathcal C_{s'}}\mathbf B\cdot d\mathbf l
    =2\pi s'B_\phi(s')=\mu_0 I
    \quad\Longrightarrow\quad
    \mathbf B(s')=\frac{\mu_0 I}{2\pi s'}\hat{\boldsymbol\phi}.
\]
Para la superficie de la espira en el plano $sz$, elegimos el
vector normal $+\hat{\boldsymbol{\phi}}$. El elemento de área
vectorial es
\begin{equation}
    d\mathbf a = \hat{\boldsymbol{\phi}}\,dz'\,ds'.
    \label{ind:area-espira}
\end{equation}
Por la regla de la mano derecha, el sentido positivo de recorrido alrededor del contorno rectangular (ubicado en $\phi = 0$) va hacia $+\hat{\mathbf{z}}$ en el borde interno $s(t)$, hacia $+\hat{\mathbf{s}}$ en el borde superior $z=b$, hacia $-\hat{\mathbf{z}}$ en el borde externo $s(t)+a$, y hacia $-\hat{\mathbf{s}}$ en el borde inferior $z=0$.

\begin{EDfigura}
\begin{tikzpicture}[
  >=Stealth,font=\normalsize,line cap=round,line join=round,
  cuerpo/.style={line width=.85pt},
  vector/.style={-{Stealth[length=2.2mm]},line width=.68pt},
  medida/.style={{Stealth[length=1.55mm]}-{Stealth[length=1.55mm]},line width=.38pt},
  ayuda/.style={black!65,line width=.38pt}
]
  \draw[ayuda,->] (0,0)--(8.40,0) node[right] {$s$};
  \draw[ayuda,->] (0,-.38)--(0,4.10) node[above] {$z$};
  \draw[cuerpo] (0,-.35)--(0,3.82);
  \node[anchor=west] at (.14,3.34) {Alambre};
  \draw[vector] (0,1.14)--(0,2.37) node[midway,left=4pt] {$I$};

  \fill[black!13] (4.63,0) rectangle (5.10,2.75);
  \draw[ayuda,dashed] (4.63,0)--(4.63,2.75) (5.10,0)--(5.10,2.75);
  \draw[cuerpo] (2.55,0) rectangle (7.00,2.75);
  \draw[vector] (2.55,.65)--(2.55,1.62) node[midway,left=4pt] {$I_{\mathrm{ind}}$};
  \draw[vector] (3.22,2.75)--(4.02,2.75);
  \draw[vector] (7.00,1.98)--(7.00,1.13);
  \draw[vector] (6.30,0)--(5.47,0);
  \node[left=3pt] at (2.55,2.18) {$\mathcal C_1$};
  \node[above=2pt] at (3.34,2.75) {$\mathcal C_2$};
  \node[right=2pt] at (7.00,2.18) {$\mathcal C_3$};
  \node[below=2pt] at (5.63,0) {$\mathcal C_4$};
  \draw[vector] (6.23,3.94)--(7.61,3.94) node[right] {$\mathbf v$};

  \foreach \xx in {3.32,3.95,5.82,6.46}{
    \node[inner sep=1pt] at (\xx,1.47) {$\otimes$};
  }
  \node at (3.68,.84) {$\mathbf B$};
  \node[anchor=south east] (inner) at (4.52,3.08) {$s'$};
  \node[anchor=south west] (outer) at (5.20,3.08) {$s'+ds'$};
  \draw[ayuda] (inner.south east)--(4.63,2.75);
  \draw[ayuda] (outer.south west)--(5.10,2.75);

  \draw[ayuda] (0,-.41)--(0,-.72) (2.55,-.09)--(2.55,-.72)
    (7,-.09)--(7,-.72);
  \draw[medida] (0,-.63)--node[below=2pt] {$s(t)$} (2.55,-.63);
  \draw[medida] (2.55,-.63)--node[below=2pt] {$a$} (7,-.63);
  \draw[ayuda] (7.10,2.75)--(7.77,2.75);
  \draw[medida] (7.67,0)--node[right=2pt] {$b$} (7.67,2.75);
\end{tikzpicture}
\EDpie{fig:espira-flujo}{Espira vista en el plano $sz$. Las cruces indican el campo del alambre entrando en la página, en la dirección $+\hat{\boldsymbol{\phi}}$ de la normal elegida. La franja sombreada va de $s'$ a $s'+ds'$ y tiene área $b\,ds'$. A lo largo de su altura, la distancia al alambre es la misma y el campo no varía con $z$. Los tramos $\mathcal C_1,\ldots,\mathcal C_4$ siguen la orientación positiva de la espira.}
\end{EDfigura}

En la definición del flujo magnético sustituimos el campo del
alambre y el elemento de área de \eqref{ind:area-espira}. Como
$\hat{\boldsymbol{\phi}}\cdot\hat{\boldsymbol{\phi}}=1$, la
integral en $z'$ suma el flujo a lo largo de toda la altura
de la franja sombreada en la figura~\ref{fig:espira-flujo}.
\begin{align}
    \Phi_B(t)
    &=\int_{S(t)}\mathbf B\cdot d\mathbf a \notag\\
    &=\int_{s(t)}^{s(t)+a}\int_0^b
      \left(\frac{\mu_0 I}{2\pi s'}\hat{\boldsymbol{\phi}}\right)
      \cdot\left(\hat{\boldsymbol{\phi}}\,dz'\,ds'\right) \notag\\
    &=\int_{s(t)}^{s(t)+a}\frac{\mu_0 I}{2\pi s'}
      \underbrace{\left(\int_0^b dz'\right)}_{b}\,ds' \notag\\
    &=\frac{\mu_0 I b}{2\pi}
      \ln\!\left(\frac{s(t)+a}{s(t)}\right).
    \label{ind:flujo-espira}
\end{align}
Por la ley de Faraday, $\mathcal E=-d\Phi_B/dt$.
Usamos el flujo de \eqref{ind:flujo-espira}.
El campo del alambre es estacionario y el flujo cambia porque
la espira se aleja. Escribiendo $s=s(t)$ y usando $ds/dt=v_0$,
obtenemos
\begin{align}
    \mathcal{E} = -\frac{d\Phi_B}{dt} &= -\frac{\mu_0 I b}{2\pi} \frac{d}{ds}\Big[ \ln(s + a) - \ln s \Big] \frac{ds}{dt} \nonumber \\
    &= -\frac{\mu_0 I b v_0}{2\pi} \left( \frac{1}{s + a} - \frac{1}{s} \right) = \frac{\mu_0 I a b v_0}{2\pi s(s + a)}
\end{align}
Por la ley de Ohm, $\mathcal E=RI_{\text{ind}}$ en el modelo
resistivo de la espira. Como $\mathcal E>0$,
\begin{equation}
    I_{\text{ind}} = \frac{\mathcal{E}}{R} = \frac{\mu_0 I a b v_0}{2\pi R s(s + a)}
\end{equation}
y circula en el sentido positivo definido: sube ($+\hat{\mathbf{z}}$) por el lado cercano $s$ y baja ($-\hat{\mathbf{z}}$) por el lado lejano $s+a$. Esto concuerda con la Ley de Lenz: al alejarse la espira, disminuye el flujo en $+\hat{\boldsymbol{\phi}}$, por lo que la corriente inducida genera un campo propio a favor de $+\hat{\boldsymbol{\phi}}$ para oponerse a dicha disminución.

\paragraph{b) Fuerza electromotriz de movimiento:}
La fuerza de Lorentz debida al movimiento de la espira es
$\mathbf F_{\text{mag}}=q\mathbf v\times\mathbf B$. Al dividir
por $q$ y sustituir $\mathbf v=v_0\hat{\mathbf s}$ y el campo
del alambre, obtenemos
\[
    \mathbf f_{\text{mag}}
    =\mathbf v\times\mathbf B
    =\frac{\mu_0 I v_0}{2\pi s'}
       \bigl(\hat{\mathbf s}\times\hat{\boldsymbol\phi}\bigr)
    =\frac{\mu_0 I v_0}{2\pi s'}\hat{\mathbf z}.
\]
La fem de movimiento es la circulación de $\mathbf f_{\mathrm{mag}}$.
Llamamos $\mathcal E_i$ al aporte del tramo $\mathcal C_i$ de la
figura~\ref{fig:espira-flujo}.
\begin{align*}
    \mathcal E
    &=\oint_{\mathcal C}\mathbf f_{\mathrm{mag}}\cdot d\mathbf l\\
    &=\sum_{i=1}^4
      \underbrace{\int_{\mathcal C_i}\mathbf f_{\mathrm{mag}}
        \cdot d\mathbf l}_{\mathcal E_i}.
\end{align*}
En $\mathcal C_1$ y $\mathcal C_3$, usamos
$d\mathbf l=\hat{\mathbf z}\,dz'$. En $\mathcal C_2$ y $\mathcal C_4$,
usamos $d\mathbf l=\hat{\mathbf s}\,ds'$. El orden de los límites
sigue el recorrido indicado en la figura.
\begin{align*}
    \mathcal E_1
    &=\int_0^b\frac{\mu_0 I v_0}{2\pi s}
      \underbrace{(\hat{\mathbf z}\cdot\hat{\mathbf z})}_{1}\,dz'\\
    &=\frac{\mu_0 I v_0}{2\pi s}\bigl[z'\bigr]_0^b\\
    &=\frac{\mu_0 I b v_0}{2\pi s},\\[3pt]
    \mathcal E_2
    &=\int_s^{s+a}\frac{\mu_0 I v_0}{2\pi s'}
      \underbrace{(\hat{\mathbf z}\cdot\hat{\mathbf s})}_{0}\,ds'\\
    &=0,\\[3pt]
    \mathcal E_3
    &=\int_b^0\frac{\mu_0 I v_0}{2\pi(s+a)}
      \underbrace{(\hat{\mathbf z}\cdot\hat{\mathbf z})}_{1}\,dz'\\
    &=\frac{\mu_0 I v_0}{2\pi(s+a)}\bigl[z'\bigr]_b^0\\
    &=-\frac{\mu_0 I b v_0}{2\pi(s+a)},\\[3pt]
    \mathcal E_4
    &=\int_{s+a}^s\frac{\mu_0 I v_0}{2\pi s'}
      \underbrace{(\hat{\mathbf z}\cdot\hat{\mathbf s})}_{0}\,ds'\\
    &=0.
\end{align*}
Sumamos los aportes de la figura en el mismo orden,
\begin{align}
    \mathcal E
    &=\mathcal E_1+\mathcal E_2+\mathcal E_3+\mathcal E_4
      \notag\\
    &=\frac{\mu_0 I b v_0}{2\pi s}
      +0-\frac{\mu_0 I b v_0}{2\pi(s+a)}+0
      \notag\\
    &=\frac{\mu_0 I b v_0}{2\pi}
       \left(\frac{1}{s}-\frac{1}{s+a}\right)
      \notag\\
    &=\frac{\mu_0 I a b v_0}{2\pi s(s+a)}.
    \label{ind:fem-movimiento}
\end{align}
Se recupera de forma idéntica la fem obtenida mediante la derivada temporal del flujo.
\paragraph{c) Fuerza de frenado y balance energético:}

La fuerza sobre cada tramo es
$\mathbf F_i=I_{\text{ind}}\int_{\mathcal C_i}d\mathbf l\times\mathbf B$.
Usamos la misma numeración del contorno.

\begin{EDfigura}
\begin{tikzpicture}[
  >=Stealth,font=\normalsize,line cap=round,line join=round,
  cuerpo/.style={line width=.85pt},
  corriente/.style={-{Stealth[length=1.9mm]},line width=.62pt},
  fuerza/.style={-{Stealth[length=2.7mm]},line width=1.12pt}
]
  \draw[cuerpo] (0,-.45)--(0,3.55);
  \node[above] at (0,3.55) {Alambre};
  \draw[corriente] (0,.64)--(0,1.74) node[midway,left=4pt] {$I$};
  \draw[cuerpo] (3.00,0) rectangle (7.51,2.58);
  \draw[corriente] (3,.51)--(3,1.15) node[midway,right=5pt] {$I_{\mathrm{ind}}$};
  \draw[corriente] (4.06,2.58)--(4.65,2.58);
  \draw[corriente] (7.51,2.04)--(7.51,1.46);
  \draw[corriente] (6.45,0)--(5.86,0);
  \node[right=3pt] at (3,2.17) {$\mathcal C_1$};
  \node[above=2pt] at (4.37,2.58) {$\mathcal C_2$};
  \node[right=2pt] at (7.51,2.08) {$\mathcal C_3$};
  \node[below=2pt] at (6.15,0) {$\mathcal C_4$};

  \draw[fuerza] (3.00,1.70)--(.83,1.70)
    node[midway,above=4pt] {$\mathbf F_{\mathrm{cercano}}$};
  \draw[fuerza] (7.51,1.15)--(8.70,1.15)
    node[midway,below=4pt] {$\mathbf F_{\mathrm{lejano}}$};
  \draw[fuerza] (5.25,2.58)--(5.25,3.42)
    node[above=2pt] {$\mathbf F_{\mathrm{sup}}$};
  \draw[fuerza] (5.25,0)--(5.25,-.84)
    node[below=2pt] {$\mathbf F_{\mathrm{inf}}$};
  \draw[corriente] (7.49,3.30)--(8.71,3.30) node[right] {$\mathbf v$};
  \foreach \xx in {4.11,4.87,5.63,6.39}{
    \node at (\xx,1.20) {$\otimes$};
  }
  \node at (5.25,.64) {$\mathbf B$};
\end{tikzpicture}
\EDpie{fig:espira-fuerzas}{Fuerzas magnéticas sobre los cuatro lados de la espira. Las corrientes paralelas del alambre y del lado cercano se atraen. El lado lejano es repelido con menor fuerza porque el campo allí es más débil. Las fuerzas verticales se cancelan y la fuerza neta apunta en sentido contrario a $\mathbf v$.}
\end{EDfigura}
Tomamos los mismos elementos de línea y los mismos límites que
en el inciso (b). Al sustituir en la definición de $\mathbf F_i$,
los productos cruz fijan la dirección de cada fuerza.
\begin{align*}
    \mathbf F_1
    &=I_{\mathrm{ind}}\frac{\mu_0 I}{2\pi s}
      \int_0^b(\hat{\mathbf z}\times\hat{\boldsymbol\phi})\,dz'\\
    &=-\frac{\mu_0 I I_{\mathrm{ind}}b}{2\pi s}\hat{\mathbf s},\\[3pt]
    \mathbf F_2
    &=I_{\mathrm{ind}}\frac{\mu_0 I}{2\pi}
      \int_s^{s+a}
      (\hat{\mathbf s}\times\hat{\boldsymbol\phi})\,\frac{ds'}{s'}\\
    &=+\frac{\mu_0 I I_{\mathrm{ind}}}{2\pi}
      \ln\!\left(\frac{s+a}{s}\right)\hat{\mathbf z},\\[3pt]
    \mathbf F_3
    &=I_{\mathrm{ind}}\frac{\mu_0 I}{2\pi(s+a)}
      \int_b^0(\hat{\mathbf z}\times\hat{\boldsymbol\phi})\,dz'\\
    &=+\frac{\mu_0 I I_{\mathrm{ind}}b}{2\pi(s+a)}\hat{\mathbf s},\\[3pt]
    \mathbf F_4
    &=I_{\mathrm{ind}}\frac{\mu_0 I}{2\pi}
      \int_{s+a}^s
      (\hat{\mathbf s}\times\hat{\boldsymbol\phi})\,\frac{ds'}{s'}\\
    &=-\frac{\mu_0 I I_{\mathrm{ind}}}{2\pi}
      \ln\!\left(\frac{s+a}{s}\right)\hat{\mathbf z}.
\end{align*}
En los tramos superior ($+\hat{\mathbf{s}}$) e inferior ($-\hat{\mathbf{s}}$), las fuerzas apuntan en $+\hat{\mathbf{z}}$ y $-\hat{\mathbf{z}}$ respectivamente, con idéntica magnitud en cada punto $s'$, cancelándose mutuamente.
\begin{align}
    \mathbf F_{\mathrm{mag}}
    &=\mathbf F_1+\mathbf F_2+\mathbf F_3+\mathbf F_4
      \notag\\
    &=\mathbf F_1+\mathbf F_3
      \notag\\
    &=-\frac{\mu_0 I I_{\mathrm{ind}} b}{2\pi}
      \left(\frac{1}{s}-\frac{1}{s+a}\right)\hat{\mathbf s}
      \notag\\
    &=-\frac{\mu_0 I I_{\mathrm{ind}}ab}{2\pi s(s+a)}
      \hat{\mathbf s}
      \notag\\
    &=-\frac{v_0}{R}
      \left[\frac{\mu_0 Iab}{2\pi s(s+a)}\right]^2\hat{\mathbf s}.
\end{align}
Apunta en $-\hat{\mathbf{s}}$, frenando el movimiento de alejamiento. Para sostener velocidad constante se requiere $\mathbf{F}_{\text{ext}} = -\mathbf{F}_{\text{mag}}$. La potencia mecánica suministrada por el agente externo es:
Al sustituir la fem de \eqref{ind:fem-movimiento} en la potencia,
\begin{align}
    P_{\text{mec}}
    &=\mathbf F_{\text{ext}}\cdot\mathbf v
      \notag\\
    &=\frac{1}{R}
      \left[\frac{\mu_0 Iabv_0}{2\pi s(s+a)}\right]^2
      \underbrace{(\hat{\mathbf s}\cdot\hat{\mathbf s})}_{1}
      \notag\\
    &=\frac{\mathcal E^2}{R}
      \notag\\
    &=I_{\mathrm{ind}}^2R=P_{\text{Joule}}.
\end{align}
Todo el trabajo mecánico realizado contra el frenado magnético se transforma íntegramente en calor disipado en la resistencia.

\vspace{0.4cm}
\hrule
\vspace{0.5cm}

\subsection*{Solución Problema 2}

\paragraph{a) Campo magnético cuasiestático:}

En un solenoide finito las líneas de campo salen por un extremo
y regresan por el otro. Al alargar el solenoide, el retorno se
aleja de la zona central. En el límite infinito, queda un campo
uniforme dentro y nulo fuera. La figura~\ref{fig:solenoide-finito-infinito}
ilustra ese límite, que comprobamos con la ley de Ampère.

\begin{EDfigura}
\begin{tikzpicture}[
  font=\small, line cap=round, line join=round,
  devanado/.style={draw=black!50,line width=.52pt},
  campo/.style={line width=.85pt},
  sentido/.style={campo,-{Stealth[length=2mm]}}
]
  \node at (0,1.78) {(a) Solenoide finito};
  \foreach \x in {-1.44,-1.08,...,1.44}{
    \draw[devanado] (\x,0) ellipse[x radius=.13,y radius=.63];
  }
  \foreach \lado in {-1,1}{
    \begin{scope}[yscale=\lado]
      \draw[campo] (-1.55,.24) -- (1.55,.24)
        .. controls (2.38,.24) and (2.35,1.37) .. (1.09,1.37)
        -- (-1.09,1.37)
        .. controls (-2.35,1.37) and (-2.38,.24) .. (-1.55,.24);
      \draw[sentido] (-.64,.24)--(.04,.24);
      \draw[sentido] (.64,1.37)--(-.04,1.37);
    \end{scope}
  }

  \begin{scope}[xshift=5.65cm]
    \node at (0,1.78) {(b) Solenoide infinito};
    \node at (0,1.16) {$\mathbf B_{\mathrm{ext}}=\mathbf 0$};
    \foreach \x in {-1.44,-1.08,...,1.44}{
      \draw[devanado] (\x,0) ellipse[x radius=.13,y radius=.63];
    }
    \node at (-1.95,.52) {$\cdots$};
    \node at (1.95,.52) {$\cdots$};
    \foreach \y in {-.28,0,.28}{
      \draw[campo] (-2.05,\y)--(2.05,\y);
      \draw[sentido] (-.63,\y)--(.08,\y);
    }
    \node at (0,-.96) {$\mathbf B_{\mathrm{int}}=B_z\hat{\mathbf z}$};
  \end{scope}
\end{tikzpicture}
\EDpie{fig:solenoide-finito-infinito}{En el solenoide finito las
líneas retornan por el exterior. Los puntos suspensivos indican
que el devanado ideal continúa en ambos sentidos del eje $z$.}
\end{EDfigura}

A un instante fijo, la simetría del solenoide infinito da
$\mathbf B=B_z(s,t)\hat{\mathbf z}$. Elegimos un lazo rectangular
orientado como en la figura~\ref{fig:solenoide-ampere} y aplicamos
la ley de Ampère,
\[
    \oint_{\mathcal C}\mathbf B\cdot d\mathbf l
    =\mu_0 I_{\mathrm{enc}}.
\]

\begin{EDfigura}
\begin{tikzpicture}[
  >=Stealth, font=\small, line cap=round, line join=round,
  recorrido/.style={->,line width=.75pt},
  lazo/.style={densely dashed,line width=.55pt}
]
  \foreach \dx/\izq/\caso in {0/1.95/1,5.26/.73/2}{
    \begin{scope}[xshift=\dx cm]
      % Ejes, devanado y superficie del lazo.
      \fill[black!9] (\izq,.42) rectangle (3.29,2.24);
      \draw[thin,->] (0,0)--(4.12,0) node[right] {$s$};
      \draw[thin,->] (0,0)--(0,2.57) node[above] {$z$};
      \draw[thin] (-.07,.42)--(.07,.42) (-.07,2.24)--(.07,2.24);
      \node[left=1pt] at (-.07,.42) {$0$};
      \node[left=1pt] at (-.07,2.24) {$\ell$};
      \draw[line width=.65pt] (1.43,-.12)--(1.43,2.64);
      \node[below=2pt] at (1.43,-.12) {$s=R$};
      \foreach \h in {.61,.99,1.37,1.75,2.13}{
        \node[inner sep=1pt] at (1.43,\h) {$\otimes$};
      }

      % Recorrido: sube a la izquierda y baja a la derecha.
      \draw[lazo] (\izq,.42) rectangle (3.29,2.24);
      \draw[recorrido] (\izq,.92)--(\izq,1.68);
      \draw[recorrido] (3.29,1.68)--(3.29,.92);
      \pgfmathsetmacro{\medio}{(\izq+3.29)/2}
      \draw[recorrido] ({\medio-.24},2.24)--({\medio+.24},2.24);
      \draw[recorrido] ({\medio+.24},.42)--({\medio-.24},.42);
      \node[anchor=west,inner sep=2pt] at (\izq,1.27) {$\mathcal C_1$};
      \node[below=1pt] at (\medio,2.24) {$\mathcal C_2$};
      \node[anchor=east,inner sep=2pt] at (3.29,1.27) {$\mathcal C_3$};
      \node[above=1pt] at (\medio,.42) {$\mathcal C_4$};

      \draw[thin] (3.38,.42)--(3.78,.42) (3.38,2.24)--(3.78,2.24);
      \draw[thin,<->] (3.68,.42)--node[right=2pt] {$\ell$}
        (3.68,2.24);

      \ifnum\caso=1
        \node at (2.02,3.17) {(a) Lazo exterior};
        \node[above=2pt] at (\izq,2.24) {$s_1$};
        \node[above=2pt] at (3.29,2.24) {$s_2$};
        \node at (2.02,-.85) {$I_{\mathrm{enc}}=0$};
      \else
        \node at (2.02,3.17) {(b) Lazo que cruza $s=R$};
        \node[above=2pt] at (\izq,2.24) {$s_1<R$};
        \node[above=2pt] at (3.29,2.24) {$s_2>R$};
        \node at (2.02,-.85) {$I_{\mathrm{enc}}=n\ell I(t)$};
      \fi
    \end{scope}
  }
\end{tikzpicture}
\EDpie{fig:solenoide-ampere}{Cortes en el plano $sz$. Las cruces
representan la corriente que entra en la página en la dirección
$+\hat{\boldsymbol\phi}$. En ambos lazos, $\mathcal C_1$ sube y
$\mathcal C_3$ baja. Los tramos $\mathcal C_2$ y $\mathcal C_4$
son radiales.}
\end{EDfigura}

En ambos lazos, los lados verticales tienen
$d\mathbf l=\hat{\mathbf z}\,dz'$ y los radiales,
$d\mathbf l=\hat{\mathbf s}\,ds'$. Los límites de cada integral
siguen las flechas de la figura~\ref{fig:solenoide-ampere}.
\begin{align*}
    \oint_{\mathcal C}\mathbf B\cdot d\mathbf l
    &=\sum_{i=1}^4\int_{\mathcal C_i}\mathbf B\cdot d\mathbf l.
\end{align*}
\begin{align*}
    \int_{\mathcal C_1}\mathbf B\cdot d\mathbf l
    &=\int_0^\ell B_z(s_1,t)
      \underbrace{(\hat{\mathbf z}\cdot\hat{\mathbf z})}_{1}\,dz'
      =\ell B_z(s_1,t),\\[3pt]
    \int_{\mathcal C_2}\mathbf B\cdot d\mathbf l
    &=\int_{s_1}^{s_2} B_z(s',t)
      \underbrace{(\hat{\mathbf z}\cdot\hat{\mathbf s})}_{0}\,ds'
      =0,\\[3pt]
    \int_{\mathcal C_3}\mathbf B\cdot d\mathbf l
    &=\int_\ell^0 B_z(s_2,t)
      \underbrace{(\hat{\mathbf z}\cdot\hat{\mathbf z})}_{1}\,dz'
      =-\ell B_z(s_2,t),\\[3pt]
    \int_{\mathcal C_4}\mathbf B\cdot d\mathbf l
    &=\int_{s_2}^{s_1} B_z(s',t)
      \underbrace{(\hat{\mathbf z}\cdot\hat{\mathbf s})}_{0}\,ds'
      =0.
\end{align*}
Al sumar los cuatro aportes,
\begin{equation}
    \oint_{\mathcal C}\mathbf B\cdot d\mathbf l
    =\ell\bigl[B_z(s_1,t)-B_z(s_2,t)\bigr].
    \label{ind:ampere-tramos-solenoide}
\end{equation}
En $\mathcal C_2$ y $\mathcal C_4$,
$\hat{\mathbf z}\cdot\hat{\mathbf s}=0$ porque el recorrido
radial es perpendicular al campo. En $\mathcal C_3$, el signo
menos viene de recorrer el lado desde $z'=\ell$ hasta $z'=0$.

En el panel (a), ninguna espira cruza la superficie, así que
$I_{\mathrm{enc}}=0$. Al usar \eqref{ind:ampere-tramos-solenoide}
en la ley de Ampère, queda
\[
    \ell\bigl[B_z(s_1,t)-B_z(s_2,t)\bigr]=0.
\]
Como $s_1$ y $s_2$ son arbitrarios, el campo exterior es constante.
Para el solenoide aislado, el campo se anula lejos del eje, así
que esta constante vale cero.

En el panel (b), la superficie del lazo corta $n\ell$ espiras
en una altura $\ell$. Cada espira aporta $I(t)$ en el sentido
de la normal del lazo, de modo que $I_{\mathrm{enc}}=n\ell I(t)$.
Como ahora $\mathcal C_3$ está en el exterior, su contribución
también se anula. Así, \eqref{ind:ampere-tramos-solenoide} da
\[
    \ell B_z(s_1,t)=\mu_0 n\ell I(t),\qquad s_1<R.
\]
Cancelamos $\ell$ y usamos $I(t)=I_0+\alpha t$. Resulta
\begin{equation}
    \mathbf B(s,t)=
    \begin{cases}
        \mu_0 n(I_0+\alpha t)\hat{\mathbf z}, & s<R,\\
        \mathbf 0, & s>R.
    \end{cases}
\end{equation}

\paragraph{b) Campo eléctrico inducido y verificación diferencial:}

Por simetría cilíndrica, el campo inducido es azimutal,
$\mathbf E=E_\phi(s)\hat{\boldsymbol\phi}$. Usamos la ley de
Faraday sobre una circunferencia $\mathcal C_s$ y su disco $S_s$:
\[
    \oint_{\mathcal C_s}\mathbf E\cdot d\mathbf l
    =-\frac{d}{dt}\int_{S_s}\mathbf B\cdot d\mathbf a.
\]
El recorrido es antihorario visto desde $+z$, por lo que
$d\mathbf l=s\,d\phi\,\hat{\boldsymbol\phi}$ y
$d\mathbf a=da\,\hat{\mathbf z}$. Sustituimos en la circulación,
\begin{align*}
    \oint_{\mathcal C_s}\mathbf E\cdot d\mathbf l
    &=s E_\phi(s)\int_0^{2\pi}
      (\hat{\boldsymbol\phi}\cdot\hat{\boldsymbol\phi})\,d\phi\\
    &=2\pi sE_\phi(s).
\end{align*}
\begin{EDfigura}
\begin{tikzpicture}[
  >=Stealth,font=\normalsize,line cap=round,line join=round,
  campo/.style={-{Stealth[length=2.0mm]},line width=.90pt},
  radio/.style={-{Stealth[length=1.65mm]},thin}
]
  \foreach \xx/\ss/\condicion/\area in {0/1.02/{s<R}/{\pi s^2},6.35/1.86/{s>R}/{\pi R^2}}{
    \begin{scope}[xshift=\xx cm]
      \ifdim\ss pt<1.42pt
        \fill[black!14] (0,0) circle[radius=\ss];
      \else
        \fill[black!14] (0,0) circle[radius=1.42];
      \fi
      \draw[dashed,black!65] (0,0) circle[radius=1.42];
      \draw[line width=.75pt] (0,0) circle[radius=\ss];
      % Las flechas del recorrido y del campo tienen sentidos opuestos.
      \draw[campo] (109:\ss) arc[start angle=109,end angle=151,radius=\ss];
      \node[anchor=south east,inner sep=1.5pt] at (128:{\ss+.32}) {$d\mathbf l$};
      \draw[campo] (-22:\ss) arc[start angle=-22,end angle=-64,radius=\ss];
      \node[anchor=west,inner sep=1.5pt] at (-43:{\ss+.26}) {$\mathbf E$};
      \draw[radio] (0,0)--node[pos=.54,below=3pt,inner sep=1pt] {$R$} (150:1.42);
      \draw[radio] (0,0)--node[pos=.77,above=2pt] {$s$} (\ss,0);
      \fill (0,0) circle[radius=1.1pt];
      \node at (0,-.48) {$\odot\;\mathbf B$};
      \node at (.15,-1.54) {$\mathcal C_s$};
      \node at (0,2.12) {$\condicion$};
      \node at (0,-2.10) {$A_{\mathrm{flujo}}=\area$};
    \end{scope}
  }
\end{tikzpicture}
\EDpie{fig:solenoide-flujo}{Secciones del solenoide vistas desde $+z$. La línea discontinua marca el radio $R$ y el círculo grueso es el contorno de Faraday de radio $s$, orientado en sentido antihorario. La zona sombreada es la parte del disco atravesada por $\mathbf B$. Su área crece como $\pi s^2$ dentro del solenoide y se mantiene en $\pi R^2$ para cualquier lazo exterior.}
\end{EDfigura}
El flujo resulta del área sombreada de la figura~\ref{fig:solenoide-flujo}.
En el disco, $\mathbf B\cdot d\mathbf a=B_z\,da$, con $B_z=0$
fuera del devanado. Por tanto,
\begin{align*}
    \Phi_B(s,t)
    &=\int_{S_s}\mathbf B\cdot d\mathbf a\\
    &=\mu_0 n(I_0+\alpha t)\,A_{\mathrm{flujo}},\\
    \frac{d\Phi_B}{dt}
    &=\mu_0 n\alpha A_{\mathrm{flujo}}.
\end{align*}
Con $2\pi sE_\phi=-d\Phi_B/dt$ aplicamos el área de cada caso.
\begin{itemize}
    \item \textbf{Interior ($s<R$).} Todo el disco aporta flujo.
    \begin{align}
        E_\phi(s)
        &=-\frac{\mu_0 n\alpha\,\pi s^2}{2\pi s}
          \notag\\
        \mathbf E_{\mathrm{int}}(s)
        &=-\frac{\mu_0 n\alpha s}{2}\hat{\boldsymbol\phi}.
    \end{align}
    \item \textbf{Exterior ($s>R$).} El flujo magnético encerrado se
    satura en el área transversal completa del solenoide.
    \begin{align}
        E_\phi(s)
        &=-\frac{\mu_0 n\alpha\,\pi R^2}{2\pi s}
          \notag\\
        \mathbf E_{\mathrm{ext}}(s)
        &=-\frac{\mu_0 n\alpha R^2}{2s}\hat{\boldsymbol\phi}.
    \end{align}
\end{itemize}
Para verificar $\nabla\times\mathbf E=-\partial\mathbf B/\partial t$,
usamos el rotacional del campo azimutal,
\[
    \nabla\times\mathbf E
    =\frac{1}{s}\frac{\partial(sE_\phi)}{\partial s}\hat{\mathbf z}.
\]
\begin{align*}
    s<R:\quad
    \nabla\times\mathbf E_{\mathrm{int}}
    &=\frac{1}{s}\frac{\partial}{\partial s}
      \left(-\frac{\mu_0 n\alpha s^2}{2}\right)\hat{\mathbf z}\\
    &=-\mu_0 n\alpha\hat{\mathbf z}
      =-\frac{\partial\mathbf B_{\mathrm{int}}}{\partial t},\\[3pt]
    s>R:\quad
    \nabla\times\mathbf E_{\mathrm{ext}}
    &=\frac{1}{s}\frac{\partial}{\partial s}
      \left(-\frac{\mu_0 n\alpha R^2}{2}\right)\hat{\mathbf z}\\
    &=\mathbf 0=-\frac{\partial\mathbf B_{\mathrm{ext}}}{\partial t}.
\end{align*}
\paragraph{c) Dependencia de trayectoria en el exterior:}

En la circunferencia $s = 2R$, el campo vale $\mathbf{E}(2R) = -\frac{\mu_0 n \alpha R}{4} \hat{\boldsymbol{\phi}}$, y el elemento de arco es $d\mathbf{l} = 2R\, d\phi\, \hat{\boldsymbol{\phi}}$. El punto $P_1$ corresponde a $\phi = 0$ y $P_2$ a $\phi = \pi$ (o $-\pi$).

\begin{EDfigura}
\begin{tikzpicture}[
  >=Stealth,font=\normalsize,line cap=round,line join=round,
  trayecto/.style={line width=.78pt},
  vector/.style={-{Stealth[length=2.15mm]},line width=.75pt}
]
  \fill[black!14] (0,0) circle[radius=1.15];
  \draw[thin,densely dashed] (0,0) circle[radius=1.15];
  \draw[thin,->] (-2.97,0)--(3.00,0) node[right] {$x$};
  \draw[thin,->] (0,-2.47)--(0,2.50) node[above] {$y$};
  \draw[trayecto] (2.30,0) arc[start angle=0,end angle=180,radius=2.30];
  \draw[trayecto,densely dashed] (2.30,0) arc[start angle=0,end angle=-180,radius=2.30];
  \draw[vector] (108:2.30) arc[start angle=108,end angle=146,radius=2.30];
  \draw[vector] (-108:2.30) arc[start angle=-108,end angle=-146,radius=2.30];
  \draw[vector] (119:1.83) arc[start angle=119,end angle=81,radius=1.83];
  \draw[vector] (-78:1.83) arc[start angle=-78,end angle=-116,radius=1.83];
  \node at (-.53,1.46) {$\mathbf E$};
  \node at (-.53,-1.46) {$\mathbf E$};
  \node[anchor=west] at (2.05,1.68) {$C_1:\ 0\longrightarrow\pi$};
  \node[anchor=west] at (2.05,-1.68) {$C_2:\ 0\longrightarrow-\pi$};
  \fill (2.30,0) circle[radius=1.8pt] node[above right=2pt] {$P_1$}
    node[below right=2pt] {$(2R,0,0)$};
  \fill (-2.30,0) circle[radius=1.8pt] node[above left=2pt] {$P_2$}
    node[below left=2pt] {$(-2R,0,0)$};
  \node at (.52,.44) {$\odot\;\mathbf B$};
  \node at (-.48,-.43) {$s<R$};
\end{tikzpicture}
\EDpie{fig:solenoide-caminos}{Los dos caminos parten de $P_1$ y llegan a $P_2$. El superior $C_1$ se recorre contra el campo eléctrico y el inferior $C_2$ a favor. Si se vuelve por el camino inferior, el lazo cerrado así formado bordea un disco que atraviesa el solenoide. Los puntos indican el campo magnético en $+\hat{\mathbf z}$.}
\end{EDfigura}
\begin{itemize}
    \item \textbf{Camino 1 ($y > 0$, $\phi$ de $0$ a $\pi$):}
    \begin{align*}
        -\int_{P_1, C_1}^{P_2}\mathbf E\cdot d\mathbf l
        &=-\int_0^{\pi}\left(-\frac{\mu_0 n\alpha R}{4}\right)
          (2R)(\hat{\boldsymbol\phi}\cdot\hat{\boldsymbol\phi})\,d\phi\\
        &=\frac{\mu_0 n\alpha R^2}{2}[\phi]_0^{\pi}\\
        &=+\frac{\pi\mu_0 n\alpha R^2}{2}.
    \end{align*}
    \item \textbf{Camino 2 ($y < 0$, $\phi$ de $0$ a $-\pi$):}
    \begin{align*}
        -\int_{P_1, C_2}^{P_2}\mathbf E\cdot d\mathbf l
        &=-\int_0^{-\pi}\left(-\frac{\mu_0 n\alpha R}{4}\right)
          (2R)(\hat{\boldsymbol\phi}\cdot\hat{\boldsymbol\phi})\,d\phi\\
        &=\frac{\mu_0 n\alpha R^2}{2}[\phi]_0^{-\pi}\\
        &=-\frac{\pi\mu_0 n\alpha R^2}{2}.
    \end{align*}
\end{itemize}
\textit{Interpretación física:} Aunque localmente $\nabla \times \mathbf{E} = \mathbf{0}$ en toda la región exterior $s > R$, dicha región no es simplemente conexa al excluir el interior del solenoide. La unión del Camino 1 con el inverso del Camino 2 conforma un lazo cerrado que rodea al solenoide, encerrando un cambio de flujo no nulo ($\oint \mathbf{E}\cdot d\mathbf{l} = -\pi \mu_0 n \alpha R^2 \neq 0$). Como la integral de línea entre dos puntos depende del camino, el campo eléctrico inducido es no conservativo y no admite una representación global de la forma $\mathbf{E} = -\nabla V$.

\vspace{0.4cm}
\hrule
\vspace{0.5cm}

\subsection*{Solución Problema 3}

\paragraph{a) Conveniencia del teorema de reciprocidad:}
El campo magnético generado por una espira circular fuera de su eje de simetría involucra integrales elípticas completas sin expresión elemental simple; integrar además ese campo sobre todo el primer cuadrante infinito sería sumamente engorroso. En contraste, el teorema de reciprocidad de Neumann garantiza que $M_{12} = M_{21} = M$. Si hacemos pasar una corriente $I$ por los dos tramos rectilíneos del ángulo recto, calcular su campo magnético en $(d, d, 0)$ mediante Biot-Savart es elemental y, como $a \ll d$, el flujo sobre el pequeño disco es inmediato.

\paragraph{b) Campo magnético del alambre en ángulo recto en $(d, d, 0)$:}
Por la ley de Biot--Savart, el aporte de un tramo $\mathcal C_i$ es
\begin{equation}
    \mathbf B_i(\mathbf r)
    =\frac{\mu_0 I}{4\pi}
      \int_{\mathcal C_i}
      \frac{d\mathbf l'\times\boldsymbol\eta}{|\boldsymbol\eta|^3},
    \qquad \boldsymbol\eta=\mathbf r-\mathbf r'.
    \label{ind:biot-savart}
\end{equation}
Dividimos el conductor en las dos semirrectas de la figura.

\begin{EDfigura}
\begin{tikzpicture}[
  >=Stealth,font=\normalsize,line cap=round,line join=round,
  vector/.style={-{Stealth[length=2.1mm]},line width=.72pt},
  corriente/.style={-{Stealth[length=2.4mm]},line width=.9pt}
]
  \draw[thin,->] (-.46,0)--(7.45,0) node[right] {$x$};
  \draw[thin,->] (0,-.42)--(0,4.40) node[above] {$y$};
  \draw[line width=1.05pt] (7.10,0)--(0,0)--(0,4.04);
  \draw[corriente] (6.88,0)--(6.00,0) node[midway,below=3pt] {$I$};
  \draw[corriente] (0,3.08)--(0,3.86) node[midway,left=4pt] {$I$};
  \node[above=3pt] at (6.32,0) {$\mathcal C_1$};
  \node[right=3pt] at (0,3.50) {$\mathcal C_2$};
  \coordinate (P) at (2.88,2.88);
  \coordinate (X) at (5.04,0);
  \coordinate (Y) at (0,1.15);
  \draw[thin,densely dashed] (2.88,0)--(P)--(0,2.88);
  \draw[thin] (2.74,0)--(2.74,.14)--(2.88,.14)
    (0,2.74)--(.14,2.74)--(.14,2.88);
  \node[below=3pt] at (2.88,0) {$d$};
  \node[left=4pt] at (0,2.88) {$d$};
  \node[below left=2pt] at (0,0) {$O$};

  \fill[black!12] (P) circle[radius=.49];
  \draw[line width=.70pt] (P) circle[radius=.49];
  \draw[vector,shorten >=2pt] (X)--node[pos=.50,right=3pt] {$\boldsymbol{\eta}$} (P);
  \draw[vector,dashed,shorten >=2pt] (Y)--node[pos=.47,above left=2pt] {$\boldsymbol{\eta}_2$} (P);
  \fill (X) circle[radius=1.6pt] node[below=5pt] {$(x',0)$};
  \fill (Y) circle[radius=1.6pt] node[left=5pt] {$(0,y')$};
  \fill (P) circle[radius=1.6pt];
  \draw[vector] (P)--++(82:.49) node[above=2pt] {$a$};
  \draw[thin] (3.25,3.19)--(3.64,3.63)
      node[anchor=west] {$P=(d,d)$};
  \node[anchor=west] at (4.15,2.58) {$\otimes\;\mathbf B_1,\ \mathbf B_2$};
  \node[anchor=west] at (4.15,2.03) {$\hat{\mathbf n}_c=-\hat{\mathbf z}$};
\end{tikzpicture}
\EDpie{fig:alambre-angulo}{Vista desde $+z$. La corriente entra por el semieje $+x$ y sale por el semieje $+y$. La flecha $\boldsymbol{\eta}$ va desde un elemento fuente del primer tramo hasta el centro de la espira. Los productos cruzados de ambos tramos apuntan en $-\hat{\mathbf z}$ y sus campos se suman. Esquema con $a\ll d$.}
\end{EDfigura}
\begin{itemize}
    \item \textbf{Tramo $\mathcal C_1$ (semieje $+x$):} $\mathbf{r}' = x' \hat{\mathbf{x}}$ con $x' \in [0, \infty)$, llevando corriente en dirección $-\hat{\mathbf{x}}$, por lo que $I d\mathbf{l}' = -I\, dx'\, \hat{\mathbf{x}}$ al integrar con $x'$ creciente.
    El vector hacia el punto de observación $\mathbf{r} = d\hat{\mathbf{x}} + d\hat{\mathbf{y}}$ es $\boldsymbol{\eta} = (d - x')\hat{\mathbf{x}} + d\hat{\mathbf{y}}$, con $|\boldsymbol{\eta}| = \sqrt{(x' - d)^2 + d^2}$.
    El producto cruz es $d\mathbf{l}' \times \boldsymbol{\eta} = (-dx'\, \hat{\mathbf{x}}) \times \big((d - x')\hat{\mathbf{x}} + d\hat{\mathbf{y}}\big) = -d\, dx'\, \hat{\mathbf{z}}$.
    Por Biot-Savart:
    \begin{equation}
        \mathbf{B}_1 = -\frac{\mu_0 I d}{4\pi} \hat{\mathbf{z}} \int_0^\infty \frac{dx'}{\left[ (x' - d)^2 + d^2 \right]^{3/2}}
    \end{equation}
    Con $u=x'-d$, los límites pasan a ser $-d$ e $\infty$.
    \begin{align*}
        \int_{-d}^{\infty}\frac{du}{(u^2+d^2)^{3/2}}
        &=\left[\frac{u}{d^2\sqrt{u^2+d^2}}\right]_{-d}^{\infty}\\
        &=\frac{1}{d^2}-\frac{-d}{d^2\sqrt{2d^2}}\\
        &=\frac{1}{d^2}\left(1+\frac{1}{\sqrt{2}}\right).
    \end{align*}
    Luego:
    \begin{equation}
        \mathbf{B}_1 = -\frac{\mu_0 I}{4\pi d} \left( 1 + \frac{1}{\sqrt{2}} \right) \hat{\mathbf{z}}
    \end{equation}
    \item \textbf{Tramo $\mathcal C_2$ (semieje $+y$):} $\mathbf{r}' = y' \hat{\mathbf{y}}$ con $y' \in [0, \infty)$ y corriente en $+\hat{\mathbf{y}}$, es decir, $I d\mathbf{l}' = +I\, dy'\, \hat{\mathbf{y}}$. Con $\boldsymbol{\eta} = d\hat{\mathbf{x}} + (d - y')\hat{\mathbf{y}}$, el producto cruz da $(dy'\, \hat{\mathbf{y}}) \times (d\hat{\mathbf{x}}) = -d\, dy'\, \hat{\mathbf{z}}$. Por simetría bajo intercambio de $x'$ e $y'$, la integral es idéntica y suma constructivamente en $-\hat{\mathbf{z}}$: $\mathbf{B}_2 = \mathbf{B}_1$.
\end{itemize}
El campo magnético total en el centro de la espira es:
\begin{equation}
    \mathbf{B}(d, d, 0) = \mathbf{B}_1 + \mathbf{B}_2 = -\frac{\mu_0 I}{2\pi d} \left( 1 + \frac{1}{\sqrt{2}} \right) \hat{\mathbf{z}}
\end{equation}

\paragraph{c) Inductancia mutua:}
Bajo la aproximación $a\ll d$, podemos usar el campo del centro
en todo el disco. Tomamos la normal paralela al campo para que $M>0$.
\begin{align*}
    \Phi_c
    &=\int_{S_c}\mathbf B\cdot d\mathbf a\\
    &\approx|\mathbf B(d,d,0)|
        \int_{S_c}\underbrace{
          (-\hat{\mathbf z})\cdot(-\hat{\mathbf z})}_{1}\,da\\
    &=\frac{\mu_0 I}{2\pi d}\left(1+\frac{1}{\sqrt{2}}\right)
         \pi a^2\\
    &=\frac{\mu_0 I a^2}{2d}\left(1+\frac{1}{\sqrt{2}}\right).
\end{align*}
Dividiendo por la corriente de prueba $I$, obtenemos la inductancia mutua:
\begin{equation}
    M = \frac{\Phi_c}{I} = \frac{\mu_0 a^2}{2d} \left( 1 + \frac{1}{\sqrt{2}} \right)
\end{equation}

\paragraph{d) Fuerza electromotriz inducida en el ángulo recto:}
Si $I_c(t) = I_0 \cos(\omega t)$ recorre la espira circular en sentido horario visto desde $+z$ (de modo que su vector área apunta en $-\hat{\mathbf{z}}$, consistente con el signo positivo de $M$ respecto al sentido de referencia en el ángulo recto), la fem inducida en el alambre infinito es:
\begin{align*}
    \mathcal E(t)
    &=-M\frac{d I_c(t)}{dt}\\
    &=-M I_0\frac{d}{dt}\cos(\omega t)\\
    &=M I_0\omega\sin(\omega t)\\
    &=\frac{\mu_0 a^2 I_0\omega}{2d}
      \left(1+\frac{1}{\sqrt{2}}\right)\sin(\omega t).
\end{align*}
\vspace{0.4cm}
\hrule
\vspace{0.5cm}

\subsection*{Solución Problema 4}

\paragraph{a) Campo magnético en todo el espacio:}
Aplicamos la ley de Ampère sobre una circunferencia $\mathcal C_s$.
Por simetría, $\mathbf B=B_\phi(s)\hat{\boldsymbol\phi}$ y
$d\mathbf l=s\,d\phi\,\hat{\boldsymbol\phi}$. Así,
\[
    \oint_{\mathcal C_s}\mathbf B\cdot d\mathbf l
    =\int_0^{2\pi}B_\phi(s)s\,d\phi
    =2\pi sB_\phi(s)=\mu_0 I_{\mathrm{enc}}(s).
\]

\begin{EDfigura}
\begin{tikzpicture}[
  >=Stealth,font=\normalsize,line cap=round,line join=round,
  vector/.style={-{Stealth[length=2.1mm]},line width=.74pt},
  radio/.style={-{Stealth[length=1.6mm]},thin}
]
  \fill[black!14] (0,0) circle[radius=.76];
  \draw[line width=.75pt] (0,0) circle[radius=.76];
  \draw[line width=.95pt] (0,0) circle[radius=1.94];
  \draw[thin,densely dashed] (0,0) circle[radius=1.37];
  \draw[vector] (58:1.37) arc[start angle=58,end angle=118,radius=1.37];
  \node at (-1.07,1.15) {$\mathbf B$};
  \node at (1.16,.99) {$\mathcal C_s$};
  \draw[radio] (0,0)--node[pos=.60,above=3pt] {$a$} (18:.76);
  \draw[radio] (0,0)--node[pos=.77,below=3pt] {$b$} (-20:1.94);
  \draw[radio] (0,0)--node[pos=.74,below=3pt] {$s$} (218:1.37);
  \foreach \t in {108,182,278}{
    \node[inner sep=1pt] at (\t:.49) {$\odot$};
  }
  \foreach \t in {30,95,154,209,266,322}{
    \node[fill=white,inner sep=1pt] at (\t:1.94) {$\otimes$};
  }
  \node at (0,2.27) {Corte desde $+z$};
  \node at (0,-2.37) {$\odot\ +I\qquad\otimes\ -I$};

  \begin{scope}[shift={(3.72,-1.55)}]
    \def\posa{1.10}
    \def\posb{2.78}
    \def\altura{2.57}
    \pgfmathsetmacro{\limiteb}{\altura*\posa/\posb}
    \draw[thin,->] (-.14,0)--(4.80,0) node[right] {$s$};
    \draw[thin,->] (0,-.14)--(0,3.28) node[above] {$B_\phi(s)$};
    \draw[thin,densely dashed] (\posa,0)--(\posa,\altura)
      (\posb,0)--(\posb,\limiteb);
    \draw[line width=.95pt] (0,0)--(\posa,\altura);
    \draw[line width=.95pt,domain=\posa:\posb,samples=65,variable=\t]
      plot ({\t},{\altura*\posa/\t});
    \draw[line width=.95pt] (\posb,0)--(4.60,0);
    \draw[fill=white,line width=.65pt] (\posb,\limiteb) circle[radius=2pt]
      (\posb,0) circle[radius=2pt];
    \fill (\posa,\altura) circle[radius=1.5pt];
    \node[below left=2pt] at (0,0) {$0$};
    \node[below=3pt] at (\posa,0) {$a$};
    \node[below=3pt] at (\posb,0) {$b$};
    \node at (.37,1.83) {$\propto s$};
    \node at (2.20,1.93) {$\propto1/s$};
    \node at (3.81,.37) {$0$};
  \end{scope}
\end{tikzpicture}
\EDpie{fig:coaxial}{Cable coaxial visto desde $+z$ y perfil radial esquemático del campo. Los puntos representan la corriente saliente del núcleo y las cruces la corriente entrante de retorno en la cáscara. $B_\phi$ aumenta dentro del núcleo, decrece entre conductores y salta a cero al cruzar la cáscara ideal delgada. El lazo dibujado encierra $I$ y queda entre $a$ y $b$.}
\end{EDfigura}
Calculamos $I_{\mathrm{enc}}=\int_{S_s}\mathbf J\cdot d\mathbf a$,
incluyendo la corriente de la cáscara cuando queda dentro de
$\mathcal C_s$. En el núcleo, $J_z=I/(\pi a^2)$ es uniforme.
La normal del disco apunta en $+\hat{\mathbf z}$. Para $s<a$,
\begin{align*}
    I_{\mathrm{enc}}(s)
    &=\int_{S_s}\mathbf J\cdot d\mathbf a\\
    &=\frac{I}{\pi a^2}\int_{S_s}
      (\hat{\mathbf z}\cdot\hat{\mathbf z})\,da\\
    &=I\frac{s^2}{a^2}.
\end{align*}
Cuando $a<s<b$, el lazo encierra toda la corriente del núcleo,
$I_{\mathrm{enc}}(s)=I$. Para $s>b$, también encierra el retorno
por la cáscara y $I_{\mathrm{enc}}(s)=I-I=0$. Sustituimos estos
valores en $B_\phi(s)=\mu_0 I_{\mathrm{enc}}(s)/(2\pi s)$,
\[
    \mathbf B(s)=
    \begin{cases}
      \dfrac{\mu_0 I s}{2\pi a^2}\hat{\boldsymbol\phi}, & s<a,\\[5pt]
      \dfrac{\mu_0 I}{2\pi s}\hat{\boldsymbol\phi}, & a<s<b,\\[5pt]
      \mathbf 0, & s>b.
    \end{cases}
\]
\paragraph{b) Energía magnética por unidad de longitud:}
Partimos de $U_B=\int u_B\,dV$, con $u_B=B^2/(2\mu_0)$.
En un tramo de longitud $\ell$, una cáscara de radio $s$ y
espesor $ds$ tiene volumen $dV=2\pi s\,\ell\,ds$. Así,
\[
    \frac{U_B}{\ell}
    =\frac{1}{2\mu_0}\int_0^\infty B^2(s)\,2\pi s\,ds.
\]
Como $B=0$ para $s>b$, quedan las contribuciones del núcleo
y de la región entre los conductores.
\begin{align*}
    \frac{U_{\mathrm{int}}}{\ell}
    &=\frac{1}{2\mu_0}\int_0^a
      \left(\frac{\mu_0 I s}{2\pi a^2}\right)^2 2\pi s\,ds\\
    &=\frac{\mu_0 I^2}{4\pi a^4}
      \left[\frac{s^4}{4}\right]_0^a
      =\frac{\mu_0 I^2}{16\pi},\\[3pt]
    \frac{U_{\mathrm{ext}}}{\ell}
    &=\frac{1}{2\mu_0}\int_a^b
      \left(\frac{\mu_0 I}{2\pi s}\right)^2 2\pi s\,ds\\
    &=\frac{\mu_0 I^2}{4\pi}\int_a^b\frac{ds}{s}\\
    &=\frac{\mu_0 I^2}{4\pi}\ln\!\left(\frac{b}{a}\right).
\end{align*}
Por tanto,
\[
    \frac{U_B}{\ell}
    =\frac{\mu_0 I^2}{16\pi}
     +\frac{\mu_0 I^2}{4\pi}\ln\!\left(\frac{b}{a}\right).
\]
\paragraph{c) Autoinductancia por unidad de longitud:}
De $U_B=\tfrac12 LI^2$ obtenemos, al dividir por $\ell$,
\begin{align*}
    \mathcal L
    &=\frac{2}{I^2}\frac{U_B}{\ell}\\
    &=\frac{2}{I^2}\left[
       \frac{\mu_0 I^2}{16\pi}
       +\frac{\mu_0 I^2}{4\pi}\ln\!\left(\frac{b}{a}\right)\right]\\
    &=\underbrace{\frac{\mu_0}{8\pi}}_{\mathcal L_{\mathrm{int}}}
      +\underbrace{\frac{\mu_0}{2\pi}\ln\!\left(\frac{b}{a}\right)}
       _{\mathcal L_{\mathrm{ext}}}.
\end{align*}
Notablemente, la inductancia interna $\mathcal{L}_{\text{int}} = \frac{\mu_0}{8\pi}$ depende únicamente de la permeabilidad magnética y es independiente del radio $a$ del conductor.

\paragraph{d) Flujo interno y enlace de flujo parcial.}
El flujo magnético a través de la sección radial del núcleo se
calcula con $d\mathbf a=\hat{\boldsymbol\phi}\,dz'\,ds$.
Al integrar primero a lo largo de la longitud $\ell$, obtenemos
\begin{align*}
    \Phi_{\mathrm{int}}
    &=\int_{S_{\mathrm{int}}}\mathbf B\cdot d\mathbf a\\
    &=\int_0^a\frac{\mu_0 I s}{2\pi a^2}\,\ell\,ds\\
    &=\frac{\mu_0 I\ell}{2\pi a^2}
      \left[\frac{s^2}{2}\right]_0^a\\
    &=\frac{\mu_0 I\ell}{4\pi}.
\end{align*}
Al dividir el flujo interno sin ponderar por $I\ell$, encontramos
\[
    \frac{\Phi_{\mathrm{int}}}{I\ell}=\frac{\mu_0}{4\pi}.
\]
Al usar ese cociente como inductancia por unidad de longitud,
se obtendría el doble del valor correcto. La razón física es que el concepto elemental $L = \Phi_B/I$ asume que las líneas de campo magnético rodean el $100\%$ de la corriente $I$ (conductores filiformes). Dentro del conductor macizo ($s < a$), un tubo de flujo diferencial $d\Phi_B(s) = B(s)\ell\, ds$ ubicado a radio $s$ solo enlaza la fracción de corriente interior $f(s) = I_{\text{enc}}(s)/I = s^2/a^2$. Al ponderar cada elemento diferencial de flujo por la fracción de corriente efectivamente enlazada ($d\Lambda = f(s)\, d\Phi_B$), obtenemos:
\begin{align}
    \mathcal L_{\mathrm{int}}
    &=\frac{1}{I\ell}\int d\Lambda_{\mathrm{int}}
      \notag\\
    &=\frac{1}{I\ell}\int_0^a
      \frac{s^2}{a^2}\frac{\mu_0 I s}{2\pi a^2}\,\ell\,ds
      \notag\\
    &=\frac{\mu_0}{2\pi a^4}
      \left[\frac{s^4}{4}\right]_0^a
      \notag\\
    &=\frac{\mu_0}{8\pi}.
\end{align}
Así se recupera el valor obtenido a partir de la energía magnética.
\vspace{0.4cm}
\hrule
\vspace{0.5cm}

\subsection*{Solución Problema 5}

\paragraph{a) Campo eléctrico y corriente de desplazamiento:}
Con $Q(t) = I t$ concentrada en el origen, el campo eléctrico cuasiestático es:
\begin{equation}
    \mathbf{E}(\mathbf{r}, t) = \frac{I t}{4\pi\varepsilon_0 r^2} \hat{\mathbf{r}}
\end{equation}
La densidad de corriente de desplazamiento fuera del origen ($r > 0$) resulta estacionaria y puramente radial:
\begin{align*}
    \mathbf J_d(\mathbf r)
    &=\varepsilon_0\frac{\partial\mathbf E}{\partial t}\\
    &=\varepsilon_0\frac{\partial}{\partial t}
      \left(\frac{It}{4\pi\varepsilon_0r^2}\hat{\mathbf r}\right)\\
    &=\frac{I}{4\pi r^2}\hat{\mathbf r}.
\end{align*}
Si integramos sobre una esfera cerrada de radio $r$ centrada en el origen (con vector área exterior $d\mathbf{a} = r^2 \sin\theta'\, d\theta'\, d\phi'\, \hat{\mathbf{r}}$), el alambre ingresa por el polo sur ($\theta' = \pi$, donde $\hat{\mathbf{r}} = -\hat{\mathbf{z}}$) aportando un flujo de conducción $\int \mathbf{J} \cdot d\mathbf{a} = -I$, mientras que la corriente de desplazamiento sale uniformemente por toda la esfera aportando $\oint \mathbf{J}_d \cdot d\mathbf{a} = \frac{I}{4\pi r^2}(4\pi r^2) = +I$. Así, $\oint (\mathbf{J} + \mathbf{J}_d) \cdot d\mathbf{a} = -I + I = 0$, verificando que la corriente total generalizada de Maxwell no tiene fuentes ni sumideros.

\paragraph{b) Corriente de desplazamiento a través del casquete $S_1$:}
Recorremos el lazo circular $C$ (radio $s = r\sin\theta$) en la dirección $+\hat{\boldsymbol{\phi}}$. Por la regla de la mano derecha, la normal unitaria sobre el casquete superior $S_1$ ($0 \le \theta' \le \theta$) apunta radialmente hacia afuera ($\hat{\mathbf{n}}_1 = +\hat{\mathbf{r}}$). Integrando $\mathbf{J}_d$ sobre $S_1$:

\begin{EDfigura}
\begin{tikzpicture}[
  >=Stealth,font=\normalsize,line cap=round,line join=round,
  vector/.style={-{Stealth[length=2.0mm]},line width=.7pt},
  normal/.style={-{Stealth[length=2.45mm]},line width=.92pt},
  casquete/.style={line width=.96pt}
]
  \foreach \xx/\caso in {0/1,7.05/2}{
    \begin{scope}[xshift=\xx cm]
      \def\radio{1.84}
      \def\ang{25}
      \pgfmathsetmacro{\altura}{\radio*sin(\ang)}
      \pgfmathsetmacro{\ancho}{\radio*cos(\ang)}
      \coordinate (Cleft) at (-\ancho,\altura);
      \coordinate (Cright) at (\ancho,\altura);
      \ifnum\caso=1
        \fill[black!13] (Cright)
          arc[start angle=25,end angle=155,radius=\radio]--cycle;
      \else
        \fill[black!13] (Cleft)
          arc[start angle=155,end angle=385,radius=\radio]--cycle;
      \fi
      \draw[black!65,thin,densely dashed] (0,0) circle[radius=\radio];
      \draw[thin,->] (0,-2.38)--(0,2.42) node[above=2pt] {$+z$};
      \ifnum\caso=1
        \draw[casquete] (Cright)
          arc[start angle=25,end angle=155,radius=\radio];
      \else
        \draw[casquete] (Cleft)
          arc[start angle=155,end angle=385,radius=\radio];
      \fi
      \draw[black!70,thin,dashed] (Cleft)
        arc[start angle=180,end angle=0,x radius=\ancho,y radius=.27];
      \draw[thin] (Cleft)
        arc[start angle=180,end angle=360,x radius=\ancho,y radius=.27];
      \draw[vector,domain=222:265,samples=22,variable=\t]
        plot ({\ancho*cos(\t)},{\altura+.27*sin(\t)});
      \node at (-1.03,{\altura-.52}) {$+\hat{\boldsymbol{\phi}}$};
      \node[anchor=west] at (Cright) {$\mathcal C$};
      \draw[line width=1.08pt] (0,-2.34)--(0,0);
      \draw[vector] (0,-1.65)--(0,-.85) node[midway,left=4pt] {$I$};
      \fill (0,0) circle[radius=1.7pt];
      \ifnum\caso=1
        \node[below right=2pt] at (0,0) {$Q(t)$};
        \draw[thin] (0,0)--(Cright) node[pos=.78,below=2pt] {$r$};
        \draw[thin] (90:.49) arc[start angle=90,end angle=25,radius=.49];
        \node at (49:.74) {$\theta$};
        \draw[normal] (115:\radio)--(115:2.50)
          node[above left=1pt] {$\hat{\mathbf n}_1$};
        \draw[vector] (62:1.24)--(62:2.17) node[above right=2pt] {$\mathbf J_d$};
        \node[anchor=east] at (-1.35,1.46) {$S_1$};
      \else
        \node[below left=3pt] at (0,0) {$Q(t)$};
        \draw[normal] (-36:\radio)--(-36:1.12);
        \node at (1.20,-.51) {$\hat{\mathbf n}_2$};
        \draw[vector] (222:1.25)--(222:2.07)
          node[left=2pt] {$\mathbf J_d$};
        \node[anchor=west] at (1.42,-1.46) {$S_2$};
        \fill (0,-\radio) circle[radius=1.6pt];
      \fi
    \end{scope}
  }
\end{tikzpicture}
\EDpie{fig:ampere-casquetes}{El mismo lazo $\mathcal C$, recorrido en $+\hat{\boldsymbol{\phi}}$, limita los dos casquetes. Su radio es $s=r\sin\theta$. El trazo grueso identifica el casquete elegido. La normal $\hat{\mathbf n}_1=+\hat{\mathbf r}$ sale de la esfera y $\hat{\mathbf n}_2=-\hat{\mathbf r}$ apunta hacia el centro. La corriente de desplazamiento $\mathbf J_d$ es radial saliente en ambos casos. El alambre perfora únicamente $S_2$.}
\end{EDfigura}
En $S_1$ el elemento de área es
$d\mathbf a=+\hat{\mathbf r}\,r^2\sin\theta'\,d\theta'\,d\phi'$.
Al sustituir $\mathbf J_d$ y este elemento en la definición,
\begin{align}
    I_{d,\mathrm{enc}}^{(1)}
    &=\int_{S_1}\mathbf J_d\cdot d\mathbf a
      \notag\\
    &=\int_0^{2\pi}\!\int_0^\theta
      \frac{I}{4\pi r^2}r^2\sin\theta'
      \underbrace{(\hat{\mathbf r}\cdot\hat{\mathbf r})}_{1}
      \,d\theta'\,d\phi'
      \notag\\
    &=\frac{I}{4\pi}(2\pi)\bigl[-\cos\theta'\bigr]_0^\theta
      \notag\\
    &=\frac{I}{2}(1-\cos\theta).
    \label{ind:corriente-desplazamiento-s1}
\end{align}
\paragraph{c) Campo magnético por Ampère-Maxwell y consistencia de superficies:}
Como el alambre yace en $\theta'=\pi$, no atraviesa el casquete
$S_1$ y $I_{\text{enc}}^{(1)}=0$. En la ley de Ampère-Maxwell,
\[
    \oint_{\mathcal C}\mathbf B\cdot d\mathbf l
    =\mu_0\bigl(I_{\mathrm{enc}}+I_{d,\mathrm{enc}}\bigr),
\]
sustituimos $d\mathbf l=r\sin\theta\,d\phi\,\hat{\boldsymbol\phi}$.
Como $B_\phi(r,\theta)$ es constante sobre $\mathcal C$,
\[
    \oint_{\mathcal C}\mathbf B\cdot d\mathbf l
    =B_\phi(r,\theta)\,r\sin\theta\int_0^{2\pi}d\phi
    =2\pi r\sin\theta\,B_\phi(r,\theta).
\]
Con $I_{\mathrm{enc}}^{(1)}=0$ y
\eqref{ind:corriente-desplazamiento-s1},
\begin{align}
    2\pi r\sin\theta\,B_\phi(r,\theta)
    &=\mu_0\frac I2(1-\cos\theta)
      \notag\\
    \mathbf B(r,\theta)
    &=\frac{\mu_0 I(1-\cos\theta)}{4\pi r\sin\theta}
      \hat{\boldsymbol\phi}.
    \label{ind:campo-ampere-maxwell}
\end{align}
Si ahora elegimos el casquete complementario inferior $S_2$ ($\theta \le \theta' \le \pi$) bordeado por el \textit{mismo} lazo orientado $C$, la compatibilidad con la regla de la mano derecha exige que el vector normal apunte hacia el origen ($\hat{\mathbf{n}}_2 = -\hat{\mathbf{r}}$). Evaluando ambos términos sobre $S_2$:
\begin{itemize}
    \item En $\theta' = \pi$, $\hat{\mathbf{n}}_2 = -\hat{\mathbf{r}}|_{\theta'=\pi} = +\hat{\mathbf{z}}$, paralelo a la corriente $I$ del alambre, de modo que $I_{\text{enc}}^{(2)} = +I$.
    \item El flujo de $\mathbf J_d$ con
    $\hat{\mathbf n}_2=-\hat{\mathbf r}$ es
    \begin{align*}
        I_{d,\mathrm{enc}}^{(2)}
        &=\int_{S_2}\mathbf J_d\cdot d\mathbf a\\
        &=\int_0^{2\pi}\!\int_\theta^\pi
          \frac{I}{4\pi}\sin\theta'
          \underbrace{\bigl(\hat{\mathbf r}\cdot(-\hat{\mathbf r})\bigr)}_{-1}
          \,d\theta'\,d\phi'\\
        &=-\frac I2\bigl[-\cos\theta'\bigr]_\theta^\pi\\
        &=-\frac I2(1+\cos\theta).
    \end{align*}
\end{itemize}
Sumando ambas contribuciones en $S_2$:
\begin{equation}
    I_{\text{enc}}^{(2)} + I_{d,\text{enc}}^{(2)} = I - \frac{I}{2}(1 + \cos\theta) = \frac{I}{2}(1 - \cos\theta)
\end{equation}
Se obtiene exactamente el mismo lado derecho y el mismo campo $\mathbf{B}(r, \theta)$, demostrando cómo el término de Maxwell elimina la ambigüedad en la elección de la superficie.

\paragraph{d) Comparación con Biot-Savart e interpretación física:}
Aplicamos Biot--Savart \eqref{ind:biot-savart} al alambre,
con $z'\in(-\infty,0]$ e $I d\mathbf l'=I\,dz'\,\hat{\mathbf z}$.
Para un punto con $(s,\phi,z)=(r\sin\theta,\phi,r\cos\theta)$,
$\boldsymbol\eta=s\hat{\mathbf s}+(z-z')\hat{\mathbf z}$ y
$d\mathbf l'\times\boldsymbol\eta=s\,dz'\,\hat{\boldsymbol\phi}$.
Al sustituir,
\begin{equation}
    \mathbf{B}_{\text{alambre}} = \frac{\mu_0 I s}{4\pi} \hat{\boldsymbol{\phi}} \int_{-\infty}^0 \frac{dz'}{\left[ s^2 + (z - z')^2 \right]^{3/2}}
\end{equation}
Con $u=z'-z$, los límites pasan a ser $-\infty$ y $-z$.
\begin{align*}
    \int_{-\infty}^{-z}\frac{du}{(s^2+u^2)^{3/2}}
    &=\left[\frac{u}{s^2\sqrt{s^2+u^2}}\right]_{-\infty}^{-z}\\
    &=-\frac{z}{s^2\sqrt{s^2+z^2}}+\frac{1}{s^2}\\
    &=\frac{1}{s^2}\left(1-\frac{z}{r}\right).
\end{align*}
Sustituimos en el campo del alambre y usamos $s=r\sin\theta$
y $z=r\cos\theta$.
\begin{align*}
    \mathbf B_{\mathrm{alambre}}
    &=\frac{\mu_0 I s}{4\pi}\frac{1}{s^2}
      \left(1-\frac{z}{r}\right)\hat{\boldsymbol\phi}\\
    &=\frac{\mu_0 I(1-\cos\theta)}{4\pi r\sin\theta}\,
      \hat{\boldsymbol\phi}.
\end{align*}
El resultado coincide con el de Ampère--Maxwell
en \eqref{ind:campo-ampere-maxwell}.
Por lo tanto, el campo magnético generado por las corrientes de desplazamiento esféricas es $\mathbf{B}_d = \mathbf{B} - \mathbf{B}_{\text{alambre}} = \mathbf{0}$ en todo el espacio.

\textit{Explicación física:} La densidad de corriente de desplazamiento $\mathbf{J}_d = \frac{I}{4\pi r^2} \hat{\mathbf{r}}$ posee simetría esférica perfecta e irrotacionalidad ($\nabla \times \mathbf{J}_d = \mathbf{0}$). Si se inserta $\mathbf{J}_d$ en una integral tipo Biot-Savart, para cualquier punto de observación $\mathbf{r}$, los elementos simétricos respecto al eje definido por $\mathbf{r}$ producen productos cruz $\mathbf{J}_d(\mathbf{r}') \times (\mathbf{r} - \mathbf{r}')$ opuestos que se cancelan por pares. Así, una corriente radial esféricamente simétrica no genera campo magnético propio, pero su flujo es indispensable en la ecuación de Ampère-Maxwell para mantener la consistencia matemática del rotacional de $\mathbf{B}$ cuando la corriente de conducción termina abruptamente en un punto.

\end{document}
